\documentclass[a4paper]{article}
% Español (México) con XeLaTeX: fontspec para fuentes Unicode, polyglossia
% para la división silábica y los nombres del documento en la variante mexicana.
\usepackage{fontspec}
\usepackage{polyglossia}
\setdefaultlanguage[variant=mexican]{spanish}
% Los nombres chinos van como «pinyin (original)»: xeCJK da a los caracteres
% Han su propia fuente; sin ella XeLaTeX los omite con un aviso.
% FandolSong no cubre algunos caracteres de nombres (U+73FA); WenQuanYi Zen Hei sí.
\usepackage[AutoFallBack=true]{xeCJK}
\setCJKmainfont{FandolSong-Regular.otf}
\setCJKfallbackfamilyfont{\CJKrmdefault}{WenQuanYi Zen Hei}
% polyglossia no traduce los nombres de listings.
\AtBeginDocument{\providecommand{\lstlistingname}{}\renewcommand{\lstlistingname}{Listado}%
  \providecommand{\lstlistlistingname}{}\renewcommand{\lstlistlistingname}{Índice de listados}}
\usepackage{amsmath, amssymb}
\usepackage{graphicx}
\usepackage[margin=2.5cm]{geometry}
\usepackage[most]{tcolorbox}
\usepackage{etoolbox}
\usepackage{listings}
\usepackage{hyperref}
\usepackage{booktabs}
\usepackage{subcaption}
\usepackage{float}
\usepackage{array}

\newtcolorbox{knowledgebox}[1]{enhanced, colback=blue!5!white, colframe=blue!75!black, colbacktitle=blue!75!black, coltitle=white, fonttitle=\bfseries, title={#1}, attach boxed title to top left={yshift=-2mm, xshift=2mm}, boxrule=1pt, sharp corners}
\newtcolorbox{importantbox}[1]{enhanced, colback=yellow!10!white, colframe=yellow!80!black, colbacktitle=yellow!80!black, coltitle=black, fonttitle=\bfseries, title={#1}, sharp corners}
\newtcolorbox{warningbox}[1]{enhanced, colback=red!5!white, colframe=red!75!black, colbacktitle=red!75!black, coltitle=white, fonttitle=\bfseries, title={#1}, sharp corners}

\lstset{
    language=Python,
    basicstyle=\ttfamily\small,
    keywordstyle=\color{blue},
    stringstyle=\color{red!60!black},
    commentstyle=\color{green!60!black},
    breaklines=true,
    frame=single,
    numbers=left,
    numberstyle=\tiny\color{gray},
    captionpos=b,
    extendedchars=true
}

\newcommand{\notetitle}{[LLM Agents SP25] Reasoning, Memory \& Planning of Language Agents}
\newcommand{\noteauthors}{Elaboradas a partir de la clase de Yu Su}
\newcommand{\notedate}{2026-04-02}
\newcommand{\videochannel}{Berkeley RDI}
\newcommand{\videopublishdate}{2025-02-18}
\newcommand{\videoduration}{1:32:39}
\newcommand{\videourl}{https://www.youtube.com/watch?v=zvI4UN2_i-w}
\newcommand{\videocoverpath}{cover.jpg}
\newcommand{\repourl}{https://github.com/hqhq1025/ai-course-notes}

% Footer with repo info
\usepackage{fancyhdr}
\pagestyle{fancy}
\fancyhf{}
\fancyfoot[C]{\thepage}
\fancyfoot[R]{\footnotesize\color{gray}\href{\repourl}{AI Course Notes}}
\renewcommand{\headrulewidth}{0pt}
\renewcommand{\footrulewidth}{0pt}

\begin{document}
\begin{titlepage}
\centering
{\Large Notas del curso\par}
\vspace{1.2cm}
{\huge\bfseries \notetitle\par}
\vspace{0.8cm}
{\large \noteauthors\par}
\vspace{0.3cm}
{\large \notedate\par}
\vspace{1.2cm}
\ifdefempty{\videocoverpath}{}{
\includegraphics[width=0.82\textwidth,height=0.45\textheight,keepaspectratio]{\videocoverpath}\par
}
\vfill
\begin{tcolorbox}[width=0.9\textwidth, colback=black!2!white, colframe=black!60, sharp corners]
\textbf{Autor o canal del video}: \videochannel\par
\textbf{Fecha de publicación}: \videopublishdate\par
\textbf{Duración del video}: \videoduration\par
\ifdefempty{\videourl}{}{\textbf{Enlace del video}: \href{\videourl}{\nolinkurl{\videourl}}\par}
\textbf{Página del proyecto}: \href{\repourl}{\nolinkurl{\repourl}}\par
\end{tcolorbox}
\end{titlepage}

\tableofcontents
\newpage

\section{Ubicación del curso y tesis central}

Esta clase está impartida por el profesor Yu Su, de Ohio State University, y su tema son tres capacidades clave de los agentes de lenguaje (Language Agents): Reasoning, Memory y Planning. La clase no trata al Agent como una forma de producto única, sino como un objeto de sistema escalable: el modelo necesita percibir, decidir, ejecutar, revisar y mejorar de manera continua en un entorno dinámico.

\begin{importantbox}{Tesis central de esta clase}
\textit{``Reasoning is for better acting.''} El razonamiento no existe para exhibir pasos en un problema de papel, sino para elegir mejores acciones en un entorno real. En correspondencia, Memory no es una acumulación de bitácoras, y Planning tampoco es una lista estática de pasos; ambos están al servicio de la tasa de éxito y la estabilidad de las tareas a largo plazo.
\end{importantbox}

Yu Su comienza revisando los sesgos narrativos comunes en el auge de los Agent de los últimos dos años: por un lado, la expectativa externa de que ``2025 es el año uno de los Agent'' es sumamente alta; por otro, los equipos, en la práctica de producción, enfrentan con frecuencia problemas como costos fuera de control, colapsos en tareas de cadena larga e inconsistencia en los criterios de evaluación. El valor de este curso está precisamente en sustituir la narrativa de eslóganes por un marco de ingeniería.

\begin{warningbox}{Alta popularidad no equivale a alta disponibilidad}
\begin{itemize}
    \item Que una demo tenga éxito no significa un éxito estable bajo la distribución real de tareas;
    \item Que la calidad de una sola inferencia sea alta no significa que la calidad de la ejecución en múltiples rondas también lo sea;
    \item Que un modelo tenga puntuaciones más altas no significa que el costo total del sistema sea aceptable.
\end{itemize}
\end{warningbox}

\subsection{Resumen de la sección}
Esta sección deja establecido el criterio de juicio para toda la clase: la investigación de Agent debe centrarse en la calidad del comportamiento del sistema, y no en las capacidades de un modelo aislado. Las discusiones posteriores sobre Reasoning, Memory y Planning giran todas en torno a ``cómo elevar la tasa de éxito en tareas reales''.
\section{Reformulación del problema desde una perspectiva agent-first}

\subsection{Del centrado en el modelo al centrado en el sistema}

El curso propone la perspectiva ``agent-first'': considerar al agente como un sistema de decisión que se ejecuta de manera continua, y no como una interfaz de preguntas y respuestas de un solo uso. Este sistema contiene, como mínimo, cinco elementos básicos: representación del estado, espacio de acciones, mecanismo de memoria, mecanismo de planeación y un bucle de retroalimentación externa. Solo cuando estos cinco elementos funcionan de manera coordinada puede el agente salvar la distancia entre las tareas cortas y las tareas largas.

\begin{knowledgebox}{Abstracción mínima del agente como sistema}
El agente de lenguaje se puede simplificar como el siguiente ciclo:
\[
\text{Observe} \rightarrow \text{Reason} \rightarrow \text{Act} \rightarrow \text{Update Memory} \rightarrow \text{Replan}
\]
Cada eslabón se puede reforzar de manera independiente, pero el desempeño del sistema depende de la calidad de las interfaces entre los eslabones y del costo de latencia.
\end{knowledgebox}

\subsection{Contexto histórico y evolución de paradigmas}

Desde los primeros prototipos de AutoGPT hasta los sistemas actuales, que combinan ReAct, uso de herramientas y autorreflexión, la evolución del paradigma no ha sido ``un algoritmo que sustituye a otro'', sino más bien ``una capacidad que se suma a otra''. El curso subraya tres vías de crecimiento: el crecimiento de la calidad del razonamiento, el crecimiento de la confiabilidad de la memoria y el crecimiento de la escala de la planeación. En conjunto, determinan el límite superior de las tareas resolubles.

\begin{table}[H]
\centering
\begin{tabular}{p{2.5cm}p{4.2cm}p{4.7cm}}
\toprule
\textbf{Etapa} & \textbf{Capacidad representativa} & \textbf{Cuello de botella principal} \\
\midrule
Prototipos tempranos de agentes & Prompt + llamadas a herramientas & Propensión a la deriva y a bucles infinitos en tareas de cadena larga \\
Etapa intermedia de ingeniería & ReAct + reflexión + recuperación & Contaminación del estado, costo elevado, ajuste de parámetros complejo \\
Etapa actual de sistematización & Acciones de razonamiento + arquitectura de memoria + replaneación & La calidad del verificador y la generalización entre escenarios siguen siendo insuficientes \\
\bottomrule
\end{tabular}
\caption{La evolución de ingeniería de los agentes de lenguaje no es un avance puntual, sino una actualización simultánea de los elementos del sistema}
\end{table}

\begin{warningbox}{Malentendido frecuente: tratar al agente como un ``chatbot que sabe llamar herramientas''}
Esta clase de definición pasa por alto la gestión del estado a largo plazo y la retroalimentación de la planeación. Un sistema de agente auténtico debe manejar de manera explícita la recuperación de fallos, el cambio de estrategia, la actualización de la memoria y el restablecimiento de objetivos; de lo contrario, su capacidad se quedará limitada a tareas cortas.
\end{warningbox}

\subsection{Resumen de la sección}

La perspectiva agent-first transforma la pregunta de ``si el modelo sabe responder'' en ``si el sistema puede completar la tarea de manera estable''. Esto ofrece un criterio de evaluación unificado para los tres grandes módulos de capacidades que siguen.
\section{Reasoning: de "saber pensar" a "saber hacer"}

\subsection{El objetivo de ingeniería de Reasoning}
Yu Su insistió repetidamente en el curso: el objetivo final de Reasoning no es exhibir la cadena de razonamiento, sino mejorar la calidad de la acción (Action Quality). Si el razonamiento no produce mejores secuencias de acciones, una tasa de error más baja o una mayor eficiencia de muestreo, entonces ese razonamiento es un gasto inútil.

\begin{figure}[H]
\centering
\includegraphics[width=0.84\textwidth]{slide-02-reasoning-action.jpg}
\caption{Diapositiva hacia el minuto 00:21:50 del video: el acoplamiento entre Reasoning y Acting es el eje del curso}
\end{figure}

\begin{importantbox}{Reasoning for Acting}
\begin{itemize}
    \item Usar el razonamiento para mejorar la calidad de la selección de acciones, no para aumentar la extensión superficial del pensamiento;
    \item El razonamiento debe poder desencadenar la corrección de la estrategia (replan) y el cambio de herramienta (tool switch);
    \item El valor del razonamiento se refleja, en última instancia, en tres indicadores: tasa de éxito, costo y robustez.
\end{itemize}
\end{importantbox}

\subsection{De CoT al razonamiento ejecutable}
El curso coloca CoT, ReAct, Self-Reflection y Meta-Reasoning en un mismo espectro continuo: en esencia, todos consisten en "inyectar pensamiento estructurado antes y después de la acción". La diferencia no está en si hay razonamiento o no, sino en cómo el razonamiento conecta el estado del entorno con la interfaz de acción.

\begin{knowledgebox}{Cuatro tipos comunes de acciones de razonamiento}
\begin{enumerate}
    \item Descomposición de tareas: dividir un objetivo complejo en subobjetivos verificables;
    \item Alineación de evidencia: corregir las hipótesis internas con salidas externas de recuperación o herramientas;
    \item Atribución de fallas: identificar el origen del error (falta de conocimiento, falla de herramienta, error de planeación);
    \item Reescritura de estrategia: reescribir el plan de acciones posteriores con base en el resultado de la atribución.
\end{enumerate}
\end{knowledgebox}

\begin{warningbox}{Más tokens de razonamiento no siempre es mejor}
Sin un verificador eficaz ni restricciones de límite, una cadena de razonamiento larga puede provocar:
\begin{itemize}
    \item Explosión de costos;
    \item Amplificación de la alucinación de confianza;
    \item Que las rutas erróneas se "expliquen de un modo que parezca más correcto".
\end{itemize}
\end{warningbox}

\subsection{Resumen de la sección}
Reasoning es el mecanismo de ganancia en la decisión del Agente. Lo que realmente importa no es "si hay razonamiento o no", sino "si el razonamiento cambió la acción hacia una mejor".
\section{Language como interfaz cognitiva unificada}

\subsection{Por qué Language sigue siendo el medio central}
El curso señala que el lenguaje tiene un triple papel: medio de codificación de estado, medio de expresión de políticas, medio de protocolo de colaboración. Esto permite que el lenguaje no solo transmita explicaciones en lenguaje natural, sino que también sea vehículo de estructuras simbólicas, fragmentos de programa y registros de decisiones intermedias.

\begin{figure}[H]
\centering
\includegraphics[width=0.84\textwidth]{slide-03-reasoning-taxonomy.jpg}
\caption{Diapositiva del video, aproximadamente en 00:37:20: el curso organiza por capas los tipos de razonamiento según su punto de aplicación y su objetivo}
\end{figure}

\begin{importantbox}{Ventajas sistémicas de la interfaz de lenguaje}
\begin{itemize}
    \item Unificación entre modalidades: el texto puede servir como capa de intercambio entre visión, código y datos estructurados;
    \item Unificación entre módulos: el planner, el memory manager y el tool router pueden compartir el mismo protocolo;
    \item Unificación entre agentes: la colaboración multi-agent puede prototiparse primero mediante un protocolo de lenguaje.
\end{itemize}
\end{importantbox}

\subsection{La fusión entre el razonamiento simbólico y el razonamiento en lenguaje}
Yu Su no contrapone symbolic reasoning y neural reasoning, sino que enfatiza una «división del trabajo fusionada»: el razonamiento en lenguaje se encarga de la compresión semántica de alta dimensión y de la exploración de políticas, mientras que el razonamiento simbólico se encarga de las restricciones precisas y del cálculo verificable. La combinación de ambos es determinante para la estabilidad de las tareas de cadena larga.

\begin{knowledgebox}{Ejemplos de estrategias de fusión}
\begin{itemize}
    \item El modelo de lenguaje genera planes candidatos y el módulo simbólico realiza la verificación de restricciones;
    \item El modelo de lenguaje realiza búsqueda heurística y el ejecutor de programas hace la verificación final;
    \item El modelo de lenguaje explica los errores y el sistema de reglas hace la intercepción de seguridad.
\end{itemize}
\end{knowledgebox}

\begin{warningbox}{Los límites de depender solo del razonamiento en lenguaje}
Cuando la tarea requiere alta precisión numérica, optimización combinatoria con restricciones fuertes o demostración formal, la cadena de lenguaje natural pura tiende a distorsionarse, y es necesario introducir una representación intermedia programática.
\end{warningbox}

\subsection{Resumen de la sección}
Language no es un «formato de entrada de prompts», sino la capa de comunicación unificada del sistema de Agent. Su valor está en conectar los distintos módulos de capacidad, no en sustituir a todos los módulos.

\section{Memory: del caché de contexto a una base de conocimiento sostenible}

\subsection{Diseño de memoria por capas}
El curso considera que Memory es el factor decisivo de la estabilidad del Agent. Sin memoria por capas, las tareas largas degeneran en «empezar de cero en cada turno». Esta clase recomienda distinguir al menos: working memory, episodic memory, semantic memory, procedural memory.

\begin{figure}[H]
\centering
\includegraphics[width=0.84\textwidth]{slide-04-memory-design.jpg}
\caption{Diapositiva del video en aproximadamente 00:52:40: el módulo de memoria no es una base vectorial única, sino un sistema jerárquico de lectura y escritura}
\end{figure}

\begin{table}[H]
\centering
\begin{tabular}{p{2.8cm}p{3.6cm}p{4.8cm}}
\toprule
\textbf{Capa de memoria} & \textbf{Contenido típico} & \textbf{Foco de ingeniería} \\
\midrule
Working memory & Estado del paso actual, variables temporales & Actualización rápida, baja latencia, recortable \\
Episodic memory & Trayectoria de la tarea, registros de fallas, ramas de decisión & Soporta la revisión posterior y la transferencia de estrategias \\
Semantic memory & Hechos estables, conocimiento del dominio, documentación de herramientas & Precisión de recuperación y consistencia de versiones \\
Procedural memory & Flujos de trabajo reutilizables, plantillas de estrategia, patrones de llamada a herramientas & Diseño de la granularidad de abstracción y las condiciones de activación \\
\bottomrule
\end{tabular}
\caption{División del trabajo de las cuatro capas de memoria en el sistema del Agent}
\end{table}

\begin{importantbox}{La esencia de Memory es la lectura y escritura controladas, no el almacenamiento en un solo sentido}
El sistema de memoria debe responder cuatro preguntas: qué escribir, cuándo escribir, cómo leer y cuándo olvidar. Si solo se hace «escritura total», pronto aparecerán acumulación de ruido y desajustes de recuperación, que terminan por deteriorar la calidad del razonamiento y la planificación.
\end{importantbox}

\begin{knowledgebox}{La API mínima de las operaciones de Memory}
\begin{itemize}
    \item `write(event, importance, scope)`: la escritura incluye la importancia y el alcance;
    \item `retrieve(query, budget, freshness)`: la recuperación fija de manera explícita el presupuesto y la restricción de vigencia;
    \item `compress(window)`: resume y comprime según la etapa de la tarea;
    \item `evict(policy)`: ejecuta la eliminación por caducidad y la fusión de conflictos.
\end{itemize}
\end{knowledgebox}

\begin{warningbox}{La contaminación de la memoria es la causa oculta más común de las fallas en tareas largas}
Si se recuperan de manera repetida conclusiones erróneas del historial, el Agent cometerá el mismo error de forma estable. La propagación de la contaminación debe reducirse mediante la confianza en la fuente, la detección de conflictos y mecanismos de reversión.
\end{warningbox}

\subsection{Resumen de la sección}
Memory determina si el Agent puede acumular capacidad efectiva a través de turnos y de tareas. Estar organizada por capas, ser controlable y ser reversible son los tres requisitos mínimos de un sistema de memoria de alta calidad.

\section{Planning: controlador de estrategia en tareas de largo alcance}

\subsection{Los límites del papel de la planificación}
El curso define Planning como «asignar el presupuesto de acciones futuras bajo restricciones», y no como generar una lista de pendientes estática. Una planificación de alta calidad debe permitir la replanificación dinámica (dynamic replanning), porque la retroalimentación del entorno cambia continuamente la secuencia de acciones óptima.

\begin{figure}[H]
\centering
\includegraphics[width=0.84\textwidth]{slide-05-planning-loop.jpg}
\caption{Diapositiva del video, aproximadamente en 01:08:20: la planificación necesita formar un ciclo cerrado con la ejecución y la retroalimentación, en lugar de desplegarse de una sola vez}
\end{figure}

\begin{importantbox}{Los tres elementos necesarios del ciclo cerrado de planificación}
\begin{enumerate}
    \item Descomposición del objetivo: mapear el objetivo final a hitos intermedios verificables;
    \item Monitoreo de la ejecución: comparar continuamente el estado esperado con el estado real;
    \item Reescritura del plan: activar la replanificación cuando la desviación supere el umbral.
\end{enumerate}
\end{importantbox}

\subsection{La combinación de planificación y búsqueda}
El punto de vista de Yu Su es que «la planificación no sustituye a la búsqueda, ni la búsqueda sustituye a la planificación». La planificación determina la estructura del espacio de búsqueda, y la búsqueda determina el óptimo local de la acción. En la ingeniería de sistemas conviene elegir si introducir MCTS, beam search o generación heurística de candidatos según la complejidad de la tarea.

\begin{knowledgebox}{Cuándo se necesita búsqueda explícita}
\begin{itemize}
    \item Cuando las consecuencias de la acción se manifiestan con retraso y son irreversibles;
    \item Cuando las trampas de óptimo local son evidentes;
    \item Cuando la señal del evaluador de un solo paso es débil y ruidosa.
\end{itemize}
\end{knowledgebox}

\begin{warningbox}{El costo de la sobreplanificación}
Una planificación de granularidad demasiado fina ocupa una gran cantidad de contexto y de presupuesto de cómputo, y debilita la capacidad de adaptación en tiempo real durante la fase de ejecución. La granularidad de la planificación debe corresponder a la tasa de cambio del entorno.
\end{warningbox}

\subsection{Pseudocódigo del planificador}
\begin{lstlisting}[caption={Un ciclo simplificado de replanificación, que enfatiza la ``actualización activada por desviación''}]
state = observe()
plan = planner.make_plan(goal, state)

for step in range(budget):
    action = planner.next_action(plan, state)
    result = env.execute(action)
    state = update_state(state, result)
    memory.write(result)

    if planner.need_replan(state, goal):
        plan = planner.replan(goal, state, memory.retrieve(goal))
\end{lstlisting}

\subsection{Resumen de la sección}
Lo esencial de Planning no está en «pensarlo todo con la mayor completitud posible de antemano», sino en «poder corregir con rapidez durante la ejecución». La replanificación dinámica es la línea divisoria entre el éxito y el fracaso de un Agent de largo alcance.

\section{Multi-Agent: ganancia de colaboración y costo de comunicación}

\subsection{Cuándo usar múltiples agentes}
El curso no fomenta ``paralelizar por paralelizar''. Multi-agent es adecuado para escenarios donde la tarea es descomponible, los límites de los roles son claros y el protocolo de comunicación es estable. Si la tarea tiene un acoplamiento fuerte o el estado intermedio se comparte en gran medida, un Agent monolítico con memory fuerte suele ser más estable.

\begin{knowledgebox}{Origen de la ganancia en sistemas multiagente}
\begin{itemize}
    \item la especialización de roles reduce la complejidad de la estrategia del agente monolítico;
    \item la exploración en paralelo aumenta la cobertura de soluciones candidatas;
    \item la revisión cruzada reduce el riesgo de alucinación de una sola ruta.
\end{itemize}
\end{knowledgebox}

\begin{table}[H]
\centering
\begin{tabular}{p{3.3cm}p{3.8cm}p{3.9cm}}
\toprule
\textbf{Dimensión} & \textbf{Agent monolítico} & \textbf{Multi-Agent} \\
\midrule
Gestión del contexto & Simple, pero puede congestionarse & Distribuida, pero requiere protocolo de sincronización \\
Eficiencia de ejecución & Secuencial, fácil de depurar & Alto potencial de paralelismo, coordinación compleja \\
Localización de errores & La localización centralizada es más intuitiva & Es necesario distinguir errores de rol y errores de comunicación \\
Estructura de costos & Relativamente estable y predecible & Alto costo de comunicación e inferencia repetida \\
\bottomrule
\end{tabular}
\caption{Disyuntivas entre arquitectura monolítica y multiagente}
\end{table}

\begin{importantbox}{Principios de implementación de sistemas multiagente}
Primero se diseña el protocolo y después se agregan roles. El protocolo debe incluir al menos: formato de los mensajes, alcance del estado compartido, reglas de resolución de conflictos, y una estrategia de tiempo de espera y reversión.
\end{importantbox}

\begin{warningbox}{Modos de falla más comunes en sistemas multiagente}
\begin{itemize}
    \item la superposición de responsabilidades entre roles produce trabajo duplicado;
    \item la información de comunicación se infla y consume el presupuesto de contexto;
    \item la falta de un validador unificado da lugar a que «se confirmen entre sí pero el resultado global sea incorrecto».
\end{itemize}
\end{warningbox}

\subsection{Resumen de la sección}
Multi-agent no es más fuerte por defecto, sino condicionalmente más fuerte. Lo que importa es la calidad de la descomposición de la tarea y el diseño del protocolo de comunicación, no la cantidad de agents en sí.

\section{Evaluación y operación: convertir al Agent en un sistema administrable}

\subsection{La evaluación debe cubrir toda la cadena}
Yu Su advierte: evaluar solo la respuesta final oculta una gran cantidad de problemas sistémicos. La evaluación de Agents debe cubrir la ``calidad de la trayectoria'', la ``tasa de corrección de las llamadas a herramientas'', la ``calidad de activación de la replanificación'' y la ``precisión de lectura de la memoria''.

\begin{importantbox}{Criterios sugeridos para la evaluación}
\begin{itemize}
    \item Outcome metrics: tasa de éxito de la tarea, latencia de finalización, costo total;
    \item Process metrics: tasa de corrección de los pasos intermedios, tasa de acciones inválidas, tasa de reintentos repetidos;
    \item Safety metrics: tasa de llamadas fuera de los permisos, tasa de filtración de información sensible, tiempo de recuperación de errores.
\end{itemize}
\end{importantbox}

\begin{knowledgebox}{División de tareas entre la evaluación fuera de línea y en línea}
\begin{itemize}
    \item Fuera de línea: hace regresión rápida de los cambios de algoritmo y localiza problemas de módulo;
    \item En línea: observa la desviación de la distribución real y corrige los supuestos hechos fuera de línea;
    \item Ambas deben compartir un esquema de eventos unificado; de lo contrario, es difícil cerrar el ciclo.
\end{itemize}
\end{knowledgebox}

\begin{warningbox}{El riesgo de mirar solo el leaderboard}
Una puntuación alta puede provenir del ajuste a plantillas de tareas y no de una generalización real. Los sistemas de producción deben establecer estrategias de degradación y mecanismos de intervención humana para las tareas fuera de distribución.
\end{warningbox}

\subsection{Observabilidad de ingeniería y recuperación ante fallas}
Un Agent estable necesita tracing de grano fino: el fundamento de cada decisión paso a paso, la herramienta invocada, el resultado devuelto, los aciertos de memoria y los cambios de plan deben poder auditarse. De lo contrario, localizar un problema solo se puede hacer reproduciéndolo por conjetura, y la eficiencia de la iteración es muy baja.

\subsection{Resumen de la sección}
La evaluación y la operación no son componentes adicionales, sino la infraestructura que determina si el Agent puede seguir evolucionando. Sin observabilidad no hay iteración controlable.

\section{Análisis de casos: tareas de Web, código e investigación}

\subsection{Caso uno: interacción de cadena larga del Web Agent}
En escenarios de Web, el Agent suele necesitar navegar entre páginas, analizar información semiestructurada y completar decisiones de varios pasos. El punto de vista del curso es que el fracaso de este tipo de tareas normalmente no se debe a que «el modelo no entiende la página», sino a que «la sincronización de estado y la actualización del plan no son oportunas». Los problemas típicos incluyen: cambios en el DOM de la página que invalidan la localización, interrupciones del estado de inicio de sesión, palabras objetivo sinónimas que provocan clics erróneos, y valores devueltos por las herramientas que no coinciden con el estado visual.

\begin{importantbox}{La dificultad central del Web Agent es «un control de circuito cerrado sólido en un entorno de estructura débil»}
El Agent debe completar de manera continua el circuito cerrado Observe-Reason-Act bajo un estado de interfaz inestable. Si carece de una máquina de estados explícita y de verificadores a nivel de paso, los errores se propagarán entre turnos y terminarán manifestándose como fallos aleatorios.
\end{importantbox}

Una estrategia viable es dividir las operaciones de página en tres categorías: «acciones de navegación», «acciones de extracción» y «acciones de confirmación», y configurar una lógica de verificación distinta para cada tipo. Por ejemplo, las acciones de navegación se centran en los cambios de URL o título de página, las acciones de extracción se centran en la integridad de los campos, y las acciones de confirmación se centran en si se cumplen las restricciones de negocio. Esto permite reducir la probabilidad de que «un paso erróneo arruine todo el proceso».

\begin{knowledgebox}{Recomendaciones de evaluación para el Web Agent}
\begin{itemize}
    \item Registrar la tasa de éxito a nivel de acción, no solo la tasa de éxito del resultado final;
    \item Contabilizar el número de reintentos y de retrocesos para evaluar la estabilidad del plan;
    \item Agrupar por complejidad de página, para evitar que el promedio oculte los casos difíciles.
\end{itemize}
\end{knowledgebox}

\subsection{Caso dos: tareas a nivel de repositorio del Code Agent}
Las tareas de código parecen estructuradas, pero los objetivos a nivel de repositorio suelen involucrar dependencias entre archivos, convenciones implícitas y restricciones históricas. El curso enfatiza que la vía de mejora del Code Agent no es solo «una generación de código más fuerte», sino «una comprensión más fuerte del estado del repositorio + un ciclo de retroalimentación de verificación más fuerte».

\begin{warningbox}{Hacer solo parches de una ronda amplifica la deuda técnica}
Si el Agent solo persigue que las pruebas actuales pasen, e ignora la consistencia de la arquitectura, las reglas de nomenclatura y los contratos de interfaz, el éxito a corto plazo se convertirá en un costo de mantenimiento a largo plazo. El Code Agent necesita usar lint, type check, unit test e integration test en conjunto como verificador.
\end{warningbox}

A nivel de implementación, la tarea puede dividirse en cuatro etapas: «comprensión», «modificación», «verificación» y «revisión posterior», y la salida de cada etapa se escribe en la episodic memory. El valor de hacerlo así es que, cuando la tarea falla, el sistema puede ubicar si se trata de una desviación en la comprensión de los requisitos, un error de implementación o una cobertura de verificación insuficiente, y así activar una replanificación específica.

\begin{knowledgebox}{Recomendaciones de memory para el Code Agent}
\begin{itemize}
    \item semantic memory almacena las convenciones del proyecto, las interfaces públicas y las decisiones históricas;
    \item episodic memory almacena los intentos clave y las razones de fracaso de esta tarea;
    \item procedural memory almacena scripts reutilizables (como flujos de trabajo de pruebas de regresión).
\end{itemize}
\end{knowledgebox}

\subsection{Caso tres: el circuito cerrado de evidencia del Research Agent}
Las tareas de investigación a menudo se juzgan erróneamente como «puros problemas de reasoning». Dentro del marco de Yu Su, el Research Agent también necesita un memory y un planning sólidos: el primero se usa para la consolidación de evidencia y la gestión de conflictos, y el segundo para el cambio de ruta de investigación y la asignación del presupuesto experimental.

\begin{importantbox}{Lo clave del Research Agent no es «escribir como un artículo académico», sino «que la evidencia sea rastreable»}
El sistema debe vincular cada conclusión con su fuente, su nivel de confianza y su marca de tiempo. De lo contrario, cuando el corpus de recuperación se actualice o cambien los valores devueltos por las herramientas, las conclusiones antiguas seguirán contaminando el razonamiento posterior dentro de la capa de memoria.
\end{importantbox}

Una práctica útil es adoptar una estructura ternaria de «conclusión-evidencia-contraevidencia» al escribir en la memory: cada vez que el Agent presenta una conclusión, debe adjuntar evidencia de respaldo y posibles contraejemplos. Esto puede mejorar de manera significativa la eficacia de la etapa de reflexión (reflection) y reducir las «alucinaciones de tipo explicativo».

\begin{table}[H]
\centering
\begin{tabular}{p{2.7cm}p{3.9cm}p{4.4cm}}
\toprule
\textbf{Tipo de tarea} & \textbf{Capacidad prioritaria} & \textbf{Riesgo principal} \\
\midrule
Web Agent & Sincronización de estado + replanificación dinámica & La deriva de página invalida las acciones \\
Code Agent & Circuito cerrado de verificación + memoria del repositorio & Los parches localmente óptimos erosionan la calidad global \\
Research Agent & Gestión de evidencia + mecanismo de reflexión & Contaminación de la memoria por conclusiones sin fuente \\
\bottomrule
\end{tabular}
\caption{Diferencias en la prioridad de capacidades entre los tres tipos de tareas representativas}
\end{table}

\subsection{Resumen de la sección}
El punto en común de los tres tipos de tareas es que todas requieren un circuito cerrado del sistema; la diferencia radica en que el enfoque de diseño del verificador y de la estrategia de memoria es distinto. El desarrollo del Agent debe configurar las capacidades según los atributos de la tarea, en lugar de perseguir una plantilla única y genérica.
\section{Diseño experimental: convertir "efectivo" en evidencia reproducible}

\subsection{Diseño de la matriz experimental}
El curso propone dividir los experimentos con Agentes en dos grupos: "variables de capacidad" y "variables de sistema": las variables de capacidad incluyen el reasoning budget, la estrategia de memory y la granularidad de planning; las variables de sistema incluyen la latencia de las herramientas, la tasa de aciertos de caché, la estrategia de concurrencia y las reglas de reintento ante fallos. Solo cuando ambos grupos de variables se controlan al mismo tiempo, las conclusiones del experimento tienen poder explicativo.

\begin{knowledgebox}{Unidad experimental mínima reproducible}
\begin{itemize}
    \item Fijar el conjunto de tareas y la versión de los datos;
    \item Fijar la interfaz de herramientas y los límites de permisos;
    \item Fijar la versión del modelo y el parámetro de temperatura;
    \item Registrar la trayectoria completa y los estados intermedios.
\end{itemize}
\end{knowledgebox}

\begin{table}[H]
\centering
\begin{tabular}{p{3.1cm}p{3.3cm}p{2.7cm}p{2.2cm}}
\toprule
\textbf{Elemento experimental} & \textbf{Configuración de control} & \textbf{Métrica principal} & \textbf{Métrica secundaria} \\
\midrule
Reasoning budget & 2k / 6k / 12k tokens & Tasa de éxito de la tarea & Costo por tarea \\
Memory policy & Sin memoria / memoria por recuperación / memoria jerárquica & Tasa de finalización de tareas largas & Duración de la recuperación ante errores \\
Planning policy & Plan estático / replanificación dinámica & Número promedio de pasos & Tasa de acciones inválidas \\
Verification depth & Verificación final / verificación por paso / verificación de doble capa & Tasa de acierto & Latencia \\
\bottomrule
\end{tabular}
\caption{Matriz base de experimentos de ablación para Agentes}
\end{table}

\subsection{Clasificación de fallos y método de atribución}
Sin una clasificación unificada de fallos, el experimento cae en una "magia negra del ajuste de parámetros". Se recomienda dividir los fallos en cuatro categorías: fallos de percepción, fallos de razonamiento, fallos de ejecución y fallos de verificación. Cada categoría de fallo corresponde a una acción de corrección distinta, lo que evita la ruta ineficiente de "agrandar el modelo ante cualquier problema".

\begin{importantbox}{La atribución de fallos debe ser "ejecutable"}
El resultado de la atribución debe poder mapearse directamente a una acción de ingeniería:
\begin{itemize}
    \item Fallo de percepción $\rightarrow$ reforzar el analizador sintáctico y la verificación del estado de la página;
    \item Fallo de razonamiento $\rightarrow$ ajustar la plantilla de reasoning y el disparo de la reflexión;
    \item Fallo de ejecución $\rightarrow$ optimizar el encapsulado de herramientas y la estrategia de reintento;
    \item Fallo de verificación $\rightarrow$ ampliar la cobertura del verificador y el umbral de error.
\end{itemize}
\end{importantbox}

\begin{warningbox}{Evitar el "sesgo de las muestras exitosas"}
Si solo se analizan las trayectorias exitosas, el sistema juzgará mal su propia robustez. Se debe priorizar el análisis de las muestras de fallo de alto costo (costo elevado, cadena larga, no recuperable), pues estas muestras determinan el límite superior de viabilidad en producción.
\end{warningbox}

\subsection{Ritmo de iteración en producción}
El curso no recomienda desplegar cambios grandes de una sola vez. La ruta más segura es "shadow mode $\rightarrow$ despliegue gradual con tráfico reducido $\rightarrow$ despliegue total". Cada etapa debe tener condiciones de reversión claras y un monitoreo continuo del SLO a nivel de tarea.

\begin{knowledgebox}{Recomendaciones de monitoreo en producción}
\begin{itemize}
    \item Umbral de caída de la tasa de éxito y estrategia de degradación automática;
    \item Umbral de aumento súbito de costo y compuerta de presupuesto;
    \item Umbral de tasa de error de herramientas y mecanismo de corte automático.
\end{itemize}
\end{knowledgebox}

\subsection{Resumen de la sección}
El objetivo del diseño experimental no es "demostrar que el sistema es bueno", sino "ubicar por qué el sistema es bueno o malo". Reproducible, atribuible y reversible son las tres condiciones mínimas del sistema experimental de un Agente.

\section{Plano de despliegue: pila tecnológica y coordinación organizacional}

\subsection{Pila tecnológica de referencia}
Combinando los puntos de vista del curso, una pila de agentes viable puede dividirse en seis capas: capa de interfaz, capa de orquestación, capa de memoria, capa de herramientas, capa de validación y capa de observabilidad. Cada capa debe contar con una capacidad mínima de autonomía y comunicarse mediante un protocolo estable.

\begin{table}[H]
\centering
\begin{tabular}{p{2.4cm}p{4.0cm}p{4.6cm}}
\toprule
\textbf{Capa} & \textbf{Responsabilidad principal} & \textbf{Producto clave} \\
\midrule
Capa de interfaz & Recibir tareas, verificar permisos, dar formato a las salidas & Task spec, contexto de permisos \\
Capa de orquestación & Descomponer tareas, enrutar, disparar la replanificación & Action graph, plan de ejecución \\
Capa de memoria & Lectura y escritura por niveles, compresión, descarte, resolución de conflictos & Índice y resumen de memoria \\
Capa de herramientas & Protocolo unificado de herramientas, encapsulado de errores, reintentos & Catálogo de API invocables \\
Capa de validación & Validación de pasos y de resultado final, intercepción de riesgos & Reporte de validación y etiquetas de riesgo \\
Capa de observabilidad & Rastreo de trayectorias, agregación de métricas, alertas & Trace, dashboard, bitácora de auditoría \\
\bottomrule
\end{tabular}
\caption{Estructura de referencia de seis capas para el despliegue de sistemas de agentes}
\end{table}

\subsection{Modelo de coordinación organizacional}
Un problema organizacional común en los proyectos de agentes es que el grupo de modelos, el grupo de plataforma y el grupo de negocio optimizan cada uno sus propios objetivos parciales. La lección del curso es establecer la «tasa de éxito de la tarea» como métrica guía unificada y hacer que los distintos equipos compartan el mismo conjunto de datos de trayectorias y la misma clasificación de fallas.

\begin{importantbox}{Tres prácticas para la coordinación organizacional}
\begin{enumerate}
    \item Unificar el esquema de eventos para evitar criterios distintos entre grupos;
    \item Unificar el mecanismo de reversión para mantener controlado el riesgo en producción;
    \item Unificar el marco de prioridades, corrigiendo primero las fallas de mayor pérdida.
\end{enumerate}
\end{importantbox}

\subsection{Gobernanza y cumplimiento}
En los despliegues empresariales, el agente no es solo un problema técnico: también involucra la minimización de permisos, el manejo de datos sensibles, el registro de auditoría y la atribución de responsabilidades. El curso recomienda integrar la capacidad de gobernanza dentro del sistema desde el diseño, en lugar de aplicarla como un parche de cumplimiento después del lanzamiento.

\begin{warningbox}{La capacidad de gobernanza debe establecerse por adelantado}
\begin{itemize}
    \item Sin una jerarquía de permisos, la invocación de herramientas puede generar riesgos de exceso de privilegios;
    \item Sin una bitácora de auditoría, es difícil rastrear la responsabilidad después de un incidente;
    \item Sin reglas de límite estricto, una anomalía del modelo puede propagarse rápidamente.
\end{itemize}
\end{warningbox}

\begin{knowledgebox}{Lista de verificación previa al despliegue}
\begin{itemize}
    \item ¿Se cuenta con reversión a nivel de tarea y con la posibilidad de intervención humana?
    \item ¿Se cuenta con registros a nivel de paso y con trayectorias que puedan reproducirse?
    \item ¿Se cuenta con límites de presupuesto y con mecanismos de corte ante fallas?
    \item ¿Se completaron pruebas de equipo rojo para los escenarios sensibles?
\end{itemize}
\end{knowledgebox}

\subsection{Resumen de la sección}
El despliegue de agentes es una ingeniería conjunta de «tecnología + organización + gobernanza». Solo diseñando los tres elementos de manera integrada puede el sistema atravesar de forma estable la brecha entre la demostración y la producción.
\section{Referencia rápida de terminología y patrones de diseño}

\subsection{Comparación de términos clave}
Para evitar desviaciones semánticas cuando los equipos colaboran entre módulos, esta clase permite extraer un glosario mínimo. Se recomienda fijar estas definiciones de términos en la documentación del proyecto, para reducir el costo de comunicación y las desviaciones experimentales causadas por «mismo nombre, distinto significado».

\begin{table}[H]
\centering
\begin{tabular}{p{2.8cm}p{3.6cm}p{4.1cm}}
\toprule
\textbf{Término} & \textbf{Significado en el contexto del curso} & \textbf{Puntos a cuidar al implementarlo} \\
\midrule
Reasoning action & Paso de pensamiento intermedio usado para mejorar la calidad de las acciones posteriores & Debe vincularse a una condición de activación y a un límite de presupuesto \\
Reflection & Autorrevisión de la trayectoria existente con intento de corrección & Requiere una señal de fallo; de lo contrario tiende a girar en vacío \\
Replanning & Reescritura del plan debido a una desviación del entorno & Debe definirse con claridad el umbral de replanificación, para evitar oscilaciones frecuentes \\
Working memory & Contexto de corto plazo necesario para el paso actual & Controlar su capacidad, para evitar que desplace información clave \\
Episodic memory & Registro de eventos y resultados durante el proceso de la tarea & Debe incluir origen y marca de tiempo, para permitir la revisión posterior \\
Semantic memory & Contenido estable de hechos y base de conocimiento & Gestionar versiones y conflictos, para evitar la contaminación con conocimiento obsoleto \\
Procedural memory & Estrategias, procesos y plantillas reutilizables & Mantener las condiciones de activación, para evitar generalizaciones erróneas \\
Verifier & Mecanismo que verifica la corrección intermedia o final & El propio verificador también requiere evaluación continua \\
\bottomrule
\end{tabular}
\caption{Explicación de ingeniería de los términos relacionados con Reasoning, Memory y Planning}
\end{table}

\subsection{Patrones de diseño frecuentes}
La práctica del curso puede condensarse además en varios patrones de diseño de alta frecuencia. No son soluciones universales, pero reducen de manera considerable la fricción de migrar de un prototipo de investigación a un sistema de producción.

\begin{knowledgebox}{Cuatro patrones de diseño de uso común}
\begin{enumerate}
    \item \textbf{Plan-Execute-Verify}: primero planificar, después ejecutar y luego verificar; adecuado para tareas que se pueden verificar;
    \item \textbf{Think-Act-Reflect}: alternancia entre razonamiento y ejecución; adecuado para entornos dinámicos;
    \item \textbf{Retrieve-Reason-Write}: primero obtener evidencia y después razonar; adecuado para tareas intensivas en conocimiento;
    \item \textbf{Shadow-to-Production}: primero evaluación en modo sombra y después despliegue gradual; adecuado para negocios de alto riesgo.
\end{enumerate}
\end{knowledgebox}

\begin{warningbox}{Condiciones límite al reutilizar un patrón}
El efecto del mismo patrón puede ser opuesto en distintas distribuciones de tareas. Antes de ponerlo en producción debe verificarse si se cumplen las premisas del patrón: que la tarea sea resoluble, la estabilidad de las herramientas, el presupuesto de latencia y las restricciones de permisos.
\end{warningbox}

\subsection{Resumen de la sección}
Unificar la terminología y los patrones de diseño reduce la fricción de colaboración entre equipos y aumenta la comparabilidad de los experimentos y la mantenibilidad del sistema. Este es también un paso clave para que la ingeniería de Agent pase de la ``experiencia individual'' a la ``capacidad organizacional''.

\section{Desafíos abiertos y ruta de investigación}

\subsection{Los desafíos de largo plazo que plantea el curso}
En la parte final del curso, los desafíos se concentran en tres niveles: el nivel cognitivo (razonamiento más confiable), el nivel de sistema (memoria y planificación más robustas) y el nivel organizacional (evaluación y despliegue más confiables).

\begin{figure}[H]
\centering
\includegraphics[width=0.84\textwidth]{slide-06-open-challenges.jpg}
\caption{Página de la presentación en video, aproximadamente en el minuto 01:24:30: las direcciones futuras se concentran en el razonamiento confiable, la gobernanza de la memoria y la planificación de largo alcance}
\end{figure}

\begin{importantbox}{Palancas técnicas clave para los próximos dos años}
\begin{enumerate}
    \item Razonamiento verificable: conectar los resultados del reasoning con verificadores formales;
    \item Gobernanza de la memoria: resolver la propagación de contaminación, los conflictos de versión y la eliminación por caducidad;
    \item Generalización de la planificación: mejorar la capacidad de transferencia entre tareas y reducir el costo de reajuste por tarea.
\end{enumerate}
\end{importantbox}

\begin{knowledgebox}{Por qué es necesario optimizar conjuntamente «Reasoning + Memory + Planning»}
La optimización de un solo punto cae fácilmente en el efecto del eslabón más débil: si el razonamiento mejora pero la memoria se degrada, los errores se amplifican; si la planificación mejora pero el verificador es débil, el sistema se desvía más rápido. El objetivo de la optimización conjunta es una ganancia estable a nivel de sistema, no un máximo local.
\end{knowledgebox}

\begin{warningbox}{Restricciones reales en el despliegue}
Los escenarios empresariales suelen tener presupuestos de latencia estrictos, límites de permisos y requisitos de auditoría. Si un prototipo de investigación ignora estas restricciones, al migrarlo a producción se producirá un retroceso de rendimiento considerable.
\end{warningbox}

\subsection{Resumen de la sección}
La ruta que propone el curso no consiste en perseguir un único SOTA, sino en construir una pila de sistemas de Agent que sea verificable, reversible y escalable. La competitividad de largo plazo proviene de la capacidad de ingeniería de sistemas.

\section{Síntesis y ampliación}

\subsection{Resumen de los puntos clave de toda la clase}
\begin{table}[H]
\centering
\begin{tabular}{p{2.5cm}p{3.7cm}p{3.8cm}p{2.6cm}}
\toprule
\textbf{Módulo de capacidad} & \textbf{Juicio central de la clase} & \textbf{Acción de ingeniería} & \textbf{Riesgo principal} \\
\midrule
Reasoning & Para actuar mejor, no para tener una cadena de pensamiento más larga & Incorporar reflexión, replanificación y verificación de acciones & Costo de tokens y amplificación de alucinaciones \\
Memory & La memoria en capas es el núcleo de la estabilidad en tareas largas & Diseñar estrategias de lectura, escritura, compresión y eliminación & Contaminación de la memoria y desajuste en la recuperación \\
Planning & La planificación dinámica es superior a una lista estática & Monitorear las desviaciones y activar la replanificación & El exceso de planificación reduce la eficiencia \\
Multi-agent & Una ganancia condicional, no una ganancia por defecto & Definir primero el protocolo, y después agregar roles & Sobrecarga de comunicación y desviación en la colaboración \\
Evaluation & Hay que observar las métricas de proceso de toda la cadena & Establecer tracing y un schema unificado & Ver solo la puntuación final produce evaluaciones erróneas \\
\bottomrule
\end{tabular}
\caption{Marco de implementación sistemática de Reasoning, Memory y Planning}
\end{table}

\begin{importantbox}{Resumen en una frase}
La siguiente etapa de competencia entre los Language Agent no consiste en quién pueda hacer el demo más vistoso, sino en quién pueda convertir el razonamiento, la memoria y la planificación en un sistema de ingeniería verificable, operable y escalable.
\end{importantbox}

\subsection{Lecturas adicionales}
\begin{itemize}
    \item Yao et al., \textit{ReAct: Synergizing Reasoning and Acting in Language Models}
    \item Shinn et al., \textit{Reflexion: Language Agents with Verbal Reinforcement Learning}
    \item Park et al., \textit{Generative Agents: Interactive Simulacra of Human Behavior}
    \item Wang et al., \textit{Voyager: An Open-Ended Embodied Agent with LLMs}
    \item Blogs técnicos oficiales de OpenAI / Anthropic sobre Agent tool-use y evaluation
\end{itemize}

\subsection{Sugerencias para la práctica futura}
\begin{itemize}
    \item Primero completar en un Agent monolítico el circuito mínimo de ``observable + reversible'', y después introducir múltiples agentes;
    \item Primero construir un verificador a nivel de tarea, y después ampliar el reasoning budget;
    \item Exponer las operaciones de memoria mediante una API explícita, para evitar la acumulación implícita de contexto;
    \item Para escenarios de producción, establecer estrategias de relevo manual y degradación, priorizando la garantía de la controlabilidad.
\end{itemize}

\end{document}
